%
% Integration in geschlossener Form
%
\subsubsection{Integration in geschlossener Form}

%
% buchberger:
%
Computeralgebrasysteme werden heute routinemässig verwendet, um
Integrale in geschlossener Form zu berechnen.
Die gute alte Methode von Bernoulli verlangt, dass man den Nenner
einer rationalen Funktion in Linearfaktoren zerlegen muss.
Dies sind anspruchsvolle algebraische Aufgaben, die algorithmisch
gelöst werden müssen.
Die Technik der Gröbner-Basen und der Buchberger-Algorithmus
löst eine grosse Zahl algebraischer Probleme und spielt daher in
der modernen Computeralgbra eine zentrale Rolle.
\emph{Nathanael Fässler}
\index{Nathanael Fässler}%
\index{Fässler, Nathanael}%
löst damit in Kapitel~\ref{chapter:buchberger} ein bekanntes 
NP-vollständiges Problem und gewinnt daraus auch eine Aussage
über die Komplexität der Lösung von Polynomgleichungen.

%
% resultante:
%
Die Methode von Rothstein-Trager in Kapitel~\ref{chapter:ratint} war eine neue
Technik zur Bestimmung der Summanden einer Stammfunktion ohne vorgängige
Berechnung einer vollständigen Faktorisierung.
In Kapitel~\ref{chapter:intalg} wurde sie weiter ausgebaut um feststellen
zu können, ob überhaupt eine elementare Stammfunktion existiert. 
Die Anwendung war jedoch einigermassen schwerfällig, weil der 
euklidische Algorithmus für Polynome mit rationalen Funktionen als
Koeffizienten durchgeführt werden muss.
Die Resultante, die
\emph{Jakob Gierer}
\index{Jakob Gierer}%
\index{Gierer, Jakob}%
in Kapitel~\ref{chapter:resultante} einführt, ermöglicht einen direkteren
Zugang.

%
% elliptisch:
%
Kapitel~\ref{chapter:intalg} hat eine ganze Reihe neuer
Funktionen eingeführt, um nicht elementare Integrale darzustellen.
Die Darstellung war natürlich nicht erschöpfend.
Zum Beispiel wurden in Abschnitt~\ref{buch:intalg:section:wurzel}
nur Funktionen betrachtet, die rationale Funktionen von $x$ und der
Wurzel eines quadratischen Polynoms $\vartheta(x)=\!\sqrt{ax^2+bx+c}$
waren.
In diesem einfachen Fall konnte immer eine Stammfunktion gefunden
werden.
Für eine Wurzel eines Polynoms höheren Grades funktioniert
dies nicht mehr.
Solche Integrale sind nicht nur eine mathematische Spielerei, sie
treten erstmals beim Versuch auf, den Umfang einer Ellipse zu berechnen.
Die Umkehrfunktionen dieser elliptischen Integrale ergeben die Jacobi
elliptischen Funktionen, die zum Beispiel die Differentialgleichung
des mathematischen Pendels lösen können.
\emph{Baris Catan}
\index{Baris Catan}%
\index{Catan, Baris}%
gibt in Kapitel~\ref{chapter:elliptisch} einen Überblick über diese
Konstruktionen.


